\section{Signed defect cancellation and a one-sided form bound}
\label{sd:section}

\paragraph{Status and scope.}
These are author-checked component proofs; independent mathematical audit
is pending.  They sharpen the missing-bound reduction, but the required
arbitrary-data time estimate remains unproved.  Throughout this section the
equation is the original unforced Navier--Stokes equation on $\mathbb R^3$,
with viscosity $\nu>0$ and the selected Schwartz-data classical branch.
No weighted moment hypothesis is needed.  The only project inputs are the
accepted quotient regularity, evolution and weighted-dissipation identities;
the final nonendpoint continuation implication also uses
Proposition~\ref{de:quotient-clock}, whose independent audit is pending.

Write $\rho=|w|$, $V=\rho^{1/2}w$, $Q=\|w\|_3^3/3$,
$D=D_{\mathcal Q}$, and $\sigma=-\operatorname{div}w$. Set
\[
 B_{ij}=R_iR_j\sigma=\partial_iq_j,\qquad
 B^0=B+\tfrac13\sigma I,\qquad
 a_0=\frac{9C_S^2}{8}.
\]
The matrix $B$ is the Hessian multiplier from
\eqref{eq:qdc-potential}, not a pointwise Newtonian convolution.  It is
symmetric, $\operatorname{tr}B=-\sigma$, and $B^0$ is trace free.

\subsection{A cancellation before taking absolute values}

\begin{lemma}[Scalar cancellation and three signed representations]
\label{sd:cancellation}
For every solenoidal $u\in H^m(\mathbb R^3;\mathbb R^3)$, $m\ge4$,
with its cubic minimizing representative, all the following integrals
converge absolutely and
\begin{align}
 \int_{\mathbb R^3}\sigma\rho^3\,dx&=0,\label{sd:scalar}\\
 K_u=-\int q_j A_i\partial_i u_j
 &=-\int V^{\mathsf T}S(u)V
   =\int V^{\mathsf T}BV
   =\int V^{\mathsf T}B^0V.\label{sd:signed}
\end{align}
Here $A=\rho w$ and $S(u)_{ij}=(\partial_i u_j+\partial_j u_i)/2$.
In particular,
\begin{equation}\label{sd:sharp-entry}
 |K_u|\le\frac23\|\sigma\|_2\|w\|_6^3.
\end{equation}
The number $2/3$ is the constant supplied by the trace-free algebra;
no optimality assertion is made for the nonlinear class of representatives.
\end{lemma}
\begin{proof}
The accepted div--curl result gives $w\in L^3\cap L^6$ and
$\nabla w\in L^2$. Hence $w\in L^4$ by interpolation, $\rho^3\in L^2$,
and $\rho^3|\nabla w|\in L^1$. Apply the local Sobolev chain rule, first
to bounded truncations, to $F=\rho^3w$. The derivative is dominated by
$4\rho^3|\nabla w|$, so $F\in W^{1,1}(\mathbb R^3)$ and $F\in L^1$.
The identities $\operatorname{div}w=-\sigma$ and
$w\cdot\nabla\rho=\rho\sigma$, including the zero set by the Sobolev
level-set property, give
\[
 \operatorname{div}(\rho^3w)=2\rho^3\sigma.
\]
For a standard cutoff $\chi_R$ with $|\nabla\chi_R|\le C/R$,
\[
 \left|\int\chi_R\operatorname{div}F\right|
 =\left|\int F\cdot\nabla\chi_R\right|
 \le\frac CR\|F\|_1\longrightarrow0.
\]
Since $\operatorname{div}F\in L^1$, dominated convergence proves
\eqref{sd:scalar}. This argument never divides by $\rho$ at a zero.

Next substitute $q=w-u$ into the accepted expression for $K_u$.
The extra term is $\int A\cdot\nabla(|u|^2/2)=0$ because
$\operatorname{div}A=0$. To justify this test globally,
$A|u|^2\in L^1$ by $A\in L^{3/2}$ and $u\in L^6$; also
$|A||u||\nabla u|\in L^1$ by $u\in L^\infty$,
$\nabla u\in L^3$. Cutoff errors therefore vanish and the product rule
is valid after local smooth approximation. It follows that
$K_u=-\int\rho w_iw_j\partial_i u_j=-\int V^{\mathsf T}S(u)V$.
On the other hand,
\[
 \int\rho w_iw_j\partial_i w_j=\int\rho^3\sigma=0.
\]
Subtract $\nabla u$ from $\nabla w$ and use $\nabla q=B$ to obtain
the representation with $B$. Adding $\sigma I/3$ changes this integral
by $\frac13\int\rho^3\sigma=0$. All terms involving $B$ or
$\nabla w$ lie in $L^1$ because these matrices are in $L^2$ and
$|V|^2=\rho^3\in L^2$; the $S(u)$ term is integrable since
$S(u)\in L^\infty$ and $|V|^2\in L^1$.

The Fourier symbol of $\sigma\mapsto B^0$ is
$I/3-\xi\otimes\xi/|\xi|^2$. Thus Plancherel gives
\begin{equation}\label{sd:plancherel}
 \|B^0\|_2=\sqrt{2/3}\,\|\sigma\|_2.
\end{equation}
For any vector $z$,
$|z\otimes z-|z|^2I/3|_F=\sqrt{2/3}|z|^2$.
Since $B^0$ is trace free, Cauchy--Schwarz in space now gives
\[
 |K_u|\le\sqrt{2/3}\,\|B^0\|_2\|V\|_4^2
 =\frac23\|\sigma\|_2\|w\|_6^3.
\]
\end{proof}

\begin{proposition}[A smaller fourth-power coefficient]
\label{sd:improved}
On every compact classical interval,
\begin{equation}\label{sd:improved-diff}
 Q'+\frac\nu2D\le\frac{a_0^3}{2\nu^3}\|\sigma\|_2^4Q.
\end{equation}
Consequently Corollary~\ref{cor:qe-defect-gronwall} remains valid with
$\Lambda_\sigma$ replaced by
$\widehat\Lambda_\sigma(t)=a_0^3(2\nu^3)^{-1}
\int_0^t\|\sigma(s)\|_2^4\,ds$.
The previously audited inequality is still valid; this is a separately
review-pending strengthening, not a silent replacement of its constant.
\end{proposition}
\begin{proof}
The accepted weighted dissipation gives $\|w\|_9^3\le a_0D$.
Interpolation and \eqref{sd:sharp-entry} yield
\[
 |K_u|\le\frac23\|\sigma\|_2(3Q)^{1/4}(a_0D)^{3/4}.
\]
Use $\alpha x^{3/4}\le\varepsilon x+
27\alpha^4/(256\varepsilon^3)$ with $\varepsilon=\nu/2$.
The coefficient is
\[
 \frac{27}{32}\left(\frac23\right)^4 3a_0^3
 =\frac12a_0^3.
\]
Insert this into $Q'+\nu D=K_u$ and apply the same integrating-factor
argument as before. Neither the exponent four nor the missing uniform
bound on its time integral has changed.
\end{proof}

\subsection{Only the amplifying part of a quadratic form must be bounded}

For a real symmetric matrix field $M\in L^2(\mathbb R^3)$ define
\begin{equation}\label{sd:form-definition}
 b_\nu(M)=\max\left\{0,\sup_{0\ne\psi\in C_c^\infty(\mathbb R^3;\mathbb R^3)}
 \frac{\displaystyle\int\psi^{\mathsf T}M\psi
       -\frac{4\nu}{9}\int|\nabla\psi|^2}
      {\displaystyle\int|\psi|^2}\right\}.
\end{equation}
This is a variational definition. No existence of an eigenfunction,
discrete spectrum, or spectral theorem is assumed. Put
$b_\nu(t)=b_\nu(B^0(t))$ along the branch.

\begin{theorem}[One-sided form clock]
\label{sd:form-clock}
The function $b_\nu(t)$ is nonnegative, Borel and locally bounded on the
classical lifespan. It satisfies
\begin{align}
 0\le3b_\nu(t)&\le\frac{a_0^3}{2\nu^3}\|\sigma(t)\|_2^4,
       \label{sd:form-domination}\\
 Q'(t)+\frac\nu2D(t)&\le3b_\nu(t)Q(t).\label{sd:form-diff}
\end{align}
Writing $L_b(t)=3\int_0^t b_\nu(s)\,ds$, one has
\begin{equation}\label{sd:form-budget}
 Q(t)+\frac\nu2\int_0^tD(s)\,ds\le Q(0)e^{L_b(t)}.
\end{equation}
Thus a uniform finite bound on $\int_0^t b_\nu(s)\,ds$ below a finite
horizon is sufficient for continuation to that horizon and beyond, should
it coincide with a finite maximal time. It is not necessary to estimate
the full fourth-power defect integral first.
\end{theorem}
\begin{proof}
For fixed $M\in L^2$ the potential term is continuous on $H^1$:
$H^1\hookrightarrow L^4$ and
\[
 \left|\int\psi^{\mathsf T}M\psi-\int\phi^{\mathsf T}M\phi\right|
 \le\|M\|_2\|\psi-\phi\|_4(\|\psi\|_4+\|\phi\|_4).
\]
Thus the variational inequality extends by density to every $H^1$
vector field, and the supremum can be taken over a fixed countable
dense set of smooth compactly supported fields. In particular it
applies to the accepted $V\in H^1$.

For the special matrix $M=B^0$, trace-free algebra and
\eqref{sd:plancherel} imply, with $x=\|\nabla\psi\|_2^2/\|\psi\|_2^2$,
\[
 \frac{\int\psi^{\mathsf T}B^0\psi}{\|\psi\|_2^2}
 \le\frac23\|\sigma\|_2 C_S^{3/2}x^{3/4}.
\]
Here we used $\|\psi\|_4^2\le
\|\psi\|_2^{1/2}\|\psi\|_6^{3/2}$ and the vector Sobolev inequality.
The exact Young maximization, now with $\varepsilon=4\nu/9$, gives
\[
 b_\nu(B^0)\le\frac{27}{256}
 \frac{(2/3)^4 C_S^6\|\sigma\|_2^4}{(4\nu/9)^3}
 =\frac{a_0^3}{6\nu^3}\|\sigma\|_2^4.
\]
This proves finiteness and \eqref{sd:form-domination}.

For measurability, on any compact classical interval the map
$t\mapsto w(t)$ is continuous into $L^3$ and
$\|\sigma(t)\|_2\le\frac12\sqrt{Y(t)}$ is uniformly bounded.
For a smooth compactly supported scalar $\eta$,
$\langle\sigma(t),\eta\rangle=\int w(t)\cdot\nabla\eta$ is continuous.
Density and the uniform $L^2$ bound extend this continuity to every
$L^2$ test, so $\sigma(t)$ is weakly continuous into $L^2$ there.
The same holds for $B^0(t)$ by the bounded linear multiplier.
Each fixed smooth-field Rayleigh expression in
\eqref{sd:form-definition} is therefore continuous in time. Its
countable supremum, and its maximum with zero, are Borel. The displayed
upper bound and the local boundedness of $Y$ prove local boundedness.

The form inequality for $V$, \eqref{sd:signed}, and the accepted
$D\ge\frac89\|\nabla V\|_2^2$ give
\[
 K_u\le\frac{4\nu}{9}\|\nabla V\|_2^2+b_\nu(t)\|V\|_2^2
 \le\frac\nu2D+3b_\nu(t)Q.
\]
This proves \eqref{sd:form-diff}. The nondecreasing, locally absolutely
continuous function $L_b$ permits the integrating-factor argument,
including its nonnegative dissipation term, and gives
\eqref{sd:form-budget}. Bounded $L_b$ implies bounded $\int D$;
Proposition~\ref{de:quotient-clock} and the local $H^1$ blow-up
alternative give the claimed continuation. No endpoint theorem is used.
\end{proof}

\begin{corollary}[An explicit bound for the original missing integral]
\label{sd:missing-bound-consumer}
Put $E_0=\|u_0\|_2^2$, $Y_0=\|\nabla u_0\|_2^2$,
$Q_0=\mathcal Q(u_0)$, and
$A_* =8C_S^3C_9^3Q_0/(3\nu^3)$. For every $t<T_*$,
\begin{equation}\label{sd:conditional-missing-bound}
 \int_0^t\|\sigma(s)\|_2^4\,ds
 \le\frac{E_0Y_0}{32\nu}
       \exp\!\left(A_*\exp\!\left(3\int_0^t b_\nu(s)\,ds\right)\right).
\end{equation}
An input-only bound on the new form clock would therefore supply the
original missing fourth-power bound with the displayed constants.
No such arbitrary-data form-clock bound is established in this section.
\end{corollary}
\begin{proof}
Equation~\eqref{sd:form-budget} gives
$\int_0^tD\le2Q_0e^{L_b(t)}/\nu$.
Proposition~\ref{de:quotient-clock} implies
$\sup_{s\le t}Y(s)\le Y_0\exp(A_*e^{L_b(t)})$.
Use $\|\sigma(s)\|_2^4\le Y(s)^2/16$ and
$\int_0^tY(s)\,ds\le E_0/(2\nu)$.
\end{proof}

\subsection{A concrete endpoint tail certificate}

Let $f(t,x)=\lambda_{\max}(B^0(t,x))$. Since $B^0$ is trace free,
$f\ge0$ and the elementary eigenvalue inequality gives
\begin{equation}\label{sd:positive-potential}
 \|f(t)\|_2\le\sqrt{2/3}\|B^0(t)\|_2
 =\frac23\|\sigma(t)\|_2.
\end{equation}
For a nonnegative $L^2$ function $f$, define its amplitude threshold
\begin{equation}\label{sd:threshold-definition}
 \kappa_\nu(f)=\inf\left\{k\in\mathbb Q_{>0}:
 \|(f-k)_+\|_{3/2}\le\frac{4\nu}{9C_S^2}\right\}.
\end{equation}
This is an amplitude cutoff, not a Fourier cutoff.

\begin{proposition}[Tail certificate and optimal scalar truncation algebra]
\label{sd:tail}
The admissible set in \eqref{sd:threshold-definition} is nonempty,
$t\mapsto\kappa_\nu(f(t))$ is Borel, and
\begin{equation}\label{sd:tail-chain}
 0\le b_\nu(B^0(t))\le\kappa_\nu(f(t))
 \le\frac{a_0^3}{6\nu^3}\|\sigma(t)\|_2^4.
\end{equation}
In particular, a finite time integral of this amplitude threshold is a
sufficient certificate for the bound in
\eqref{sd:conditional-missing-bound}. More generally, any measurable
$k(t)\ge0$ with finite time integral and
$\|(f(t)-k(t))_+\|_{3/2}\le4\nu/(9C_S^2)$ almost everywhere suffices.
\end{proposition}
\begin{proof}
For any admissible $k$ and any $\psi\in H^1$,
\[
 \int\psi^{\mathsf T}B^0\psi\le\int f|\psi|^2
 \le k\|\psi\|_2^2+\|(f-k)_+\|_{3/2}C_S^2\|\nabla\psi\|_2^2.
\]
The definition of $b_\nu$ therefore gives $b_\nu\le k$; take the
infimum, with no assertion that the infimum is attained.

The exact pointwise inequality
\begin{equation}\label{sd:truncation}
 (x-k)_+^{3/2}\le\frac{3\sqrt3}{16}\,k^{-1/2}x^2
 \qquad(x\ge0,\ k>0)
\end{equation}
follows by maximizing $(y-1)^{3/2}/y^2$ for $y>1$; its maximum is at
$y=4$ and equals $3\sqrt3/16$. If
$\delta=4\nu/(9C_S^2)$, integration shows that every
\[
 k>\frac{27}{256}\delta^{-3}\|f\|_2^4
\]
is admissible. Rational approximation and
\eqref{sd:positive-potential} give
\[
 \kappa_\nu(f)\le\frac{27}{256}
 \left(\frac{9C_S^2}{4\nu}\right)^3\frac{16}{81}\|\sigma\|_2^4
 =\frac{a_0^3}{6\nu^3}\|\sigma\|_2^4.
\]
This also handles $f=0$ by taking $k\downarrow0$.

For completeness, the weakly continuous $L^2$ curve $B^0(t)$ is strongly
Borel measurable: expand it in a countable orthonormal basis and take the
pointwise $L^2$ limit of the finite partial sums. The largest-eigenvalue
map is Lipschitz in the Frobenius norm, so $f(t)$ is strongly Borel in
$L^2$. For every fixed $k>0$ the map
$g\mapsto\int(g-k)_+^{3/2}$ on nonnegative $L^2$ is continuous.
Indeed, its difference is bounded by a constant times
$\|g-h\|_2(\|g\|_2+\|h\|_2)^{1/2}
(|\{g>k\}|+|\{h>k\}|)^{1/4}$, by the power difference inequality,
Cauchy--Schwarz and Chebyshev. These measures are bounded on an
$L^2$-bounded set for fixed $k$. Thus each rational admissibility set is
Borel. The infimum over these countably many tests is Borel as well.
The general measurable certificate follows from the same form inequality.
\end{proof}

\paragraph{Scaling and the remaining implication.}
Under the actual Navier--Stokes scaling,
$B^0_\lambda(t,x)=\lambda^2B^0(\lambda^2t,\lambda x)$.
Changing variables in the Rayleigh quotient gives
\[
 b_\nu(B^0_\lambda(t))=\lambda^2 b_\nu(B^0(\lambda^2t)).
\]
The form clock is therefore critical. The same scaling holds for
$\kappa_\nu$, using continuity of the tail norm in the positive cutoff
and the density of the rationals. The bounds in
\eqref{sd:form-domination} and \eqref{sd:tail-chain} run from the
\emph{unknown fourth-power integral to the new clocks}, not the reverse
as an unconditional estimate. Neither energy's square-integral budget nor
the moment-derived $L_t^2L_x^{3/2}$ budget bounds these clocks. The remaining
producer is an input-only, endpoint-uniform bound on the signed form
clock (or on a sufficient amplitude certificate) from the vector evolution.
Formula~\eqref{sd:conditional-missing-bound} is its fully explicit consumer,
not its proof. No open claim has been promoted.

\paragraph{Relation to existing criteria.}
One-sided strain and mixed critical sum-space regularity criteria are
established mechanisms; see Miller~\cite{Miller2021Sum}, especially
Theorem~1.2 and the amplitude-decomposition arguments. That work controls
the positive middle eigenvalue of the velocity strain, whereas $B^0$ here
is the trace-free Hessian of the cubic representative's gradient correction.
The cited theorem is not applied to $B^0$. The present argument is written
out above, and no novelty claim is made for the form-bound or sum-space
mechanism. On an actual finite-horizon branch, a finite form clock gives
continuation and hence a finite defect integral; no separation of those
finiteness conditions on such a branch is claimed.
